% GENERATED FILE. Edit canonical lessons, then run npm run generate:books.
% canonical-source-hash: fnv1a64:89bbab5641727bd1
% canonical-lessons: ZH-C439-daotian, ZH-C439-guoyuan, ZH-C439-nongjiale, ZH-C439-yazi, ZH-C439-sanxin

\chapter{Country Life}
\label{ch:zh-country-life}
\hlchaptermodality{439}

\begin{chapteropening}
\textbf{By the end of this chapter:} \emph{I can talk about rice paddies, an orchard, a farm stay, ducks and getting away to unwind.}

Complete the last lesson of chapter 439: I can talk about rice paddies, an orchard, a farm stay, ducks and getting away to unwind.
\end{chapteropening}

\section[dàotián]{dàotián — rice paddy; paddy field}

\label{lesson:ZH-C439-daotian}



\begin{quote}
Before the new one: say the Mandarin for a part of a whole, then the Mandarin for all of it.
\end{quote}

\subsection*{You'll want to know: dàotián}
\emph{dàotián} — \textquotedblleft{}rice paddy; paddy field\textquotedblright{}.

\begin{quote}
\emph{wǒ yéye nǎinai zhù zài dàotián pángbiān} — My grandparents live next to the rice paddies.
\end{quote}

Its characters come later; for now it is a word for the ear.

\begin{cousinweb}
\emph{dàotián} is built from \emph{dào} 'rice in the field' and \emph{tián} 'a field', both inherited from Middle Chinese, \emph{dawX} and \emph{den} (Guangyun, in Baxter's transcription), which glosses the first as rice and the second, after an older book, as land that has been ploughed.
\end{cousinweb}

\begin{grammarlens}[title={What you've built}]
One more word for getting around.
\end{grammarlens}

\subsection*{Your turn}
\begin{itemize}
\raggedright
  \item \emph{Say it:} \emph{dàotián}
  \item \emph{Say it:} \emph{dàotián}, once more
  \item \emph{Say it:} \emph{wǒ yéye nǎinai zhù zài dàotián pángbiān} — My grandparents live next to the rice paddies.
  \item \emph{Recall:} say the Mandarin for a part of a whole, then the Mandarin for all of it, then say \emph{dàotián} again
\end{itemize}

\begin{tcolorbox}[breakable,colback=teal!4,colframe=teal!35!black,title={Before you move on}]
What does dàotián mean? (\textquotedblleft{}Rice paddy; paddy field\textquotedblright{}.) Say it once more.
\end{tcolorbox}
\section[guǒyuán]{guǒyuán — orchard}

\label{lesson:ZH-C439-guoyuan}



\begin{quote}
Before the new one: say the Mandarin for all of it, then the Mandarin for a rice paddy.
\end{quote}

\subsection*{You'll want to know: guǒyuán}
\emph{guǒyuán} — \textquotedblleft{}orchard\textquotedblright{}.

\begin{quote}
\emph{zhōumò wǒmen kāichē qù guǒyuán} — At the weekend we're driving out to the orchard.
\end{quote}

Its characters come later; for now it is a word for the ear.

\begin{cousinweb}
\emph{guǒyuán} is built from \zh{果} 'fruit' and \emph{yuán} 'a garden', both inherited from Middle Chinese, \emph{kwaX} and \emph{hjwon} (Guangyun, in Baxter's transcription), which glosses the first as the fruit of a tree and the second as a garden.
\end{cousinweb}

\begin{grammarlens}[title={What you've built}]
One more word for getting around.
\end{grammarlens}

\subsection*{Your turn}
\begin{itemize}
\raggedright
  \item \emph{Say it:} \emph{guǒyuán}
  \item \emph{Say it:} \emph{guǒyuán}, once more
  \item \emph{Say it:} \emph{zhōumò wǒmen kāichē qù guǒyuán} — At the weekend we're driving out to the orchard.
  \item \emph{Recall:} say the Mandarin for all of it, then the Mandarin for a rice paddy, then say \emph{guǒyuán} again
\end{itemize}

\begin{tcolorbox}[breakable,colback=teal!4,colframe=teal!35!black,title={Before you move on}]
What does guǒyuán mean? (\textquotedblleft{}Orchard\textquotedblright{}.) Say it once more.
\end{tcolorbox}
\section[nóngjiālè]{nóngjiālè — farm stay; country guesthouse}

\label{lesson:ZH-C439-nongjiale}



\begin{quote}
Before the new one: say the Mandarin for a rice paddy, then the Mandarin for an orchard.
\end{quote}

\subsection*{You'll want to know: nóngjiālè}
\emph{nóngjiālè} — \textquotedblleft{}farm stay; country guesthouse\textquotedblright{}.

\begin{quote}
\emph{wǒmen zhù zài yì jiā nóngjiālè} — We're staying at a farm guesthouse.
\end{quote}

Its characters come later; for now it is a word for the ear.

\begin{cousinweb}
\emph{nóngjiālè} is built from \emph{nóng} 'farming', \zh{家} 'a family' and \zh{乐} 'joy', all inherited from Middle Chinese, \emph{nowng}, \emph{kae} and \emph{lak} (Guangyun, in Baxter's transcription), which glosses the first as working the fields. The name for a farm that takes guests is modern.
\end{cousinweb}

\begin{grammarlens}[title={What you've built}]
One more word for getting around.
\end{grammarlens}

\subsection*{Your turn}
\begin{itemize}
\raggedright
  \item \emph{Say it:} \emph{nóngjiālè}
  \item \emph{Say it:} \emph{nóngjiālè}, once more
  \item \emph{Say it:} \emph{wǒmen zhù zài yì jiā nóngjiālè} — We're staying at a farm guesthouse.
  \item \emph{Recall:} say the Mandarin for a rice paddy, then the Mandarin for an orchard, then say \emph{nóngjiālè} again
\end{itemize}

\begin{tcolorbox}[breakable,colback=teal!4,colframe=teal!35!black,title={Before you move on}]
What does nóngjiālè mean? (\textquotedblleft{}Farm stay; country guesthouse\textquotedblright{}.) Say it once more.
\end{tcolorbox}
\section[yāzi]{yāzi — duck}

\label{lesson:ZH-C439-yazi}



\begin{quote}
Before the new one: say the Mandarin for an orchard, then the Mandarin for a farm stay.
\end{quote}

\subsection*{You'll want to know: yāzi}
\emph{yāzi} — \textquotedblleft{}duck\textquotedblright{}.

\begin{quote}
\emph{nǎinai jiā yǒu shí zhī yāzi} — Grandma keeps ten ducks.
\end{quote}

Its characters come later; for now it is a word for the ear.

\begin{cousinweb}
\emph{yāzi} is built from \emph{yā} 'a duck' and the noun ending \emph{zi}. The first is inherited from Middle Chinese \emph{'aep} (Guangyun, in Baxter's transcription), glossed as a water bird.
\end{cousinweb}

\begin{grammarlens}[title={What you've built}]
One more word for getting around.
\end{grammarlens}

\subsection*{Your turn}
\begin{itemize}
\raggedright
  \item \emph{Say it:} \emph{yāzi}
  \item \emph{Say it:} \emph{yāzi}, once more
  \item \emph{Say it:} \emph{nǎinai jiā yǒu shí zhī yāzi} — Grandma keeps ten ducks.
  \item \emph{Recall:} say the Mandarin for an orchard, then the Mandarin for a farm stay, then say \emph{yāzi} again
\end{itemize}

\begin{tcolorbox}[breakable,colback=teal!4,colframe=teal!35!black,title={Before you move on}]
What does yāzi mean? (\textquotedblleft{}Duck\textquotedblright{}.) Say it once more.
\end{tcolorbox}
\section[sànxīn]{sànxīn — to take one's mind off things; to unwind}

\label{lesson:ZH-C439-sanxin}



\begin{quote}
Before the new one: say the Mandarin for a farm stay, then the Mandarin for a duck.
\end{quote}

\subsection*{You'll want to know: sànxīn}
\emph{sànxīn} — \textquotedblleft{}to take one's mind off things; to unwind\textquotedblright{}.

\begin{quote}
\emph{lèi de shíhou wǒ xǐhuan qù nóngcūn sànxīn} — When I'm tired, I like to go to the countryside to unwind.
\end{quote}

Its characters come later; for now it is a word for the ear.

\begin{cousinweb}
\emph{sànxīn} is built from \emph{sàn} 'to scatter' and \zh{心} 'the heart', both inherited from Middle Chinese, \emph{sanH} and \emph{sim} (Guangyun, in Baxter's transcription), which glosses the first 'to break apart' and 'to spread': cares scattered from the heart.
\end{cousinweb}

\begin{grammarlens}[title={What you've built}]
One more word for getting around.
\end{grammarlens}

\subsection*{Your turn}
\begin{itemize}
\raggedright
  \item \emph{Say it:} \emph{sànxīn}
  \item \emph{Say it:} \emph{sànxīn}, once more
  \item \emph{Say it:} \emph{lèi de shíhou wǒ xǐhuan qù nóngcūn sànxīn} — When I'm tired, I like to go to the countryside to unwind.
  \item \emph{Recall:} say the Mandarin for a farm stay, then the Mandarin for a duck, then say \emph{sànxīn} again
\end{itemize}

\begin{tcolorbox}[breakable,colback=teal!4,colframe=teal!35!black,title={Before you move on}]
What does sànxīn mean? (\textquotedblleft{}To take one's mind off things; to unwind\textquotedblright{}.) Say it once more.
\end{tcolorbox}
